% !PRINTED 83
\weberpage{83}
%<<CONT
\weq{15}{\vartheta_{00}^2\,\vartheta_{00}^2(v)\,\frac{d^2\log\vartheta_{00}(v)}{dv^2} = \vartheta_{00}\vartheta_{00}''\vartheta_{00}^2(v) + \vartheta_{11}'^2\vartheta_{11}^2(v),}
\weq{16}{\vartheta_{01}^2\,\vartheta_{01}^2(v)\,\frac{d^2\log\vartheta_{01}(v)}{dv^2} = \vartheta_{01}\vartheta_{01}''\vartheta_{01}^2(v) - \vartheta_{11}'^2\vartheta_{11}^2(v),}
\weq{17}{\vartheta_{10}^2\,\vartheta_{10}^2(v)\,\frac{d^2\log\vartheta_{10}(v)}{dv^2} = \vartheta_{10}\vartheta_{10}''\vartheta_{10}^2(v) - \vartheta_{11}'^2\vartheta_{11}^2(v),}
\weq{18}{\vartheta_{00}^2\,\vartheta_{11}^2(v)\,\frac{d^2\log\vartheta_{11}(v)}{dv^2} = \vartheta_{00}\vartheta_{00}''\vartheta_{11}^2(v) - \vartheta_{11}'^2\vartheta_{00}^2(v).}

From these, further relations can still be derived by continued differentiation. We adduce one more of these formulae, which is obtained if one differentiates (15) twice more with respect to $v$ and then sets $v=0$. If the differentiations with respect to $v$ are expressed by those with respect to $\omega$ by means of the partial differential equation (4), it follows that:
\[
\frac{d^2\log\vartheta_{00}}{d\omega^2} - 2\left(\frac{d\log\vartheta_{00}}{d\omega}\right)^2 = -\frac{\pi^2}{8}\,\vartheta_{10}^4\vartheta_{01}^4,
\]
or
\weq{19}{\frac{d}{d\omega}\left(\frac{1}{\vartheta_{00}^4}\,\frac{d\vartheta_{00}^2}{d\omega}\right) = -\frac{\pi^2}{4}\,\frac{\vartheta_{10}^4\vartheta_{01}^4}{\vartheta_{00}^2}.}

\begin{center}
\textbf{\S\,24. \gesperrt{Representation of the $\vartheta$-Functions by Infinite}}\\
\textbf{\gesperrt{Products.}}
\end{center}

The theory of the $\vartheta$-functions has now been carried far enough for their representation to follow very easily. With this the existence of these functions is then established, and the considerations made so far first receive their secure foundation. Two ways to this goal stand open to us. The first proceeds from the zero points of the $\vartheta$-functions already known to us and composes infinite products out of them.

The function $\vartheta_{00}(v)$ vanishes, according to \S\,21, when
\[
2v = (2\nu-1)\omega + (2\mu-1)
\]
holds, where $\nu, \mu$ are integers, or when
\[
e^{\pm 2\pi i v} = -e^{\pi i \omega(2\nu-1)},
\]
where we can now restrict $\nu$ to \gesperrt{positive} numbers. If we therefore set, for brevity,
\[
e^{\pi i \omega} = q,
\]
%CONT>>
